\documentclass[11pt]{article}
\usepackage[margin=1in]{geometry}
\usepackage{amsmath, amssymb}
\usepackage{graphicx}

% augmented matrix: \begin{amatrix}{3} ... \end{amatrix} puts a bar after column 3
\newenvironment{amatrix}[1]{\left[\begin{array}{*{#1}{r}|r}}{\end{array}\right]}

\begin{document}

\textbf{DSC 210, Fall 2026 --- Homework 2} \hfill Nicholas Wendt (200478549)

%%%%%%%%%%%%%%%%%%%%%%%%%%%%%%%%%%%%%%%%%%%%%%%%%%%%%%%%%%%%%%%%%%%%%%%%
\section*{Question 1: Gauss elimination (20 points)}

Use Gauss elimination to solve the following equations for $\mathbf{x} = \begin{bmatrix} x_1 & x_2 & x_3 \end{bmatrix}^\top$:
\begin{align*}
    x_1 + 5x_2 + 3x_3 &= 1\\
    2x_1 + 4x_2 - 3x_3 &= -1\\
    -2x_1 + 2x_2 + 3x_3 &= -5
\end{align*}

\textbf{Solution:} First set up an augmented matrix
$$
\begin{amatrix}{3}
    1&5&3&1\\
    2&4&-3&-1\\
    -2&2&3&-5
\end{amatrix}
$$
Now solve the system $Ax=b$ by row reducing the matrix.\\
Apply the row operations $2r_1-r_2 \rightarrow r_2$ and $2r_1+r_3 \rightarrow r_3$
$$
\begin{amatrix}{3}
    1&5&3&1\\
    0&6&9&3\\
    0&12&9&-3
\end{amatrix}
$$
Now use row 2's pivot column $2r_2-r_3 \rightarrow r_3$
$$
\begin{amatrix}{3}
    1&5&3&1\\
    0&6&9&3\\
    0&0&9&9
\end{amatrix}
$$
Now use back substitution to solve for $\mathbf{x}$
Starting with the last row we get the equation $9x_3=9$ so $x_3=1$.\\
Now the second row gives us the equation $6x_2+9x_3 = 3 \rightarrow 6x_2+9=3$. So we see $x_2 = -1$.\\
Finally, the first equation gives us $1x_1+5x_2+3x_3 = 1 \rightarrow x_1-5+3=1$. So we see $x_1 = 3$.
$$
\mathbf{x} = \begin{bmatrix}
    x_1&x_2&x_3
\end{bmatrix}^\top=\begin{bmatrix}
    3&-1&1
\end{bmatrix}
$$
\bigskip

%%%%%%%%%%%%%%%%%%%%%%%%%%%%%%%%%%%%%%%%%%%%%%%%%%%%%%%%%%%%%%%%%%%%%%%%
\newpage
\section*{Question 2: LU decomposition (20 points)}

Perform LU decomposition for the matrix corresponding to the equations in \textbf{Question 1}. Using the matrices $L$ and $U$, do forward and backward substitution to solve for $\mathbf{x}$. Match your solution with that of \textbf{Question 1}.

\textbf{Solution:}
$$
A=\begin{bmatrix}
    1&5&3\\
    2&4&-3\\
    -2&2&3
\end{bmatrix}
$$
We can factor $A$ into a lower triangular matrix ($L$) and an upper triagular matrix ($U$) such that $A=LU$\\
We row reduce $A$ to an upper triangular matrix without using row interchanges.\\
We keep track of the scalers in $L$.\\

\bigskip
First 2 operations $r_2-2r_1$ , $r_3+2r_1$. We keep track of the scalers of the pivot row in $L$ according to
$r_i-cr_j$ where $c$ is the row $i$ column $j$ entry in $L$.\\

Now we have:
$$
A=\begin{bmatrix}
    1&5&3\\
    0&-6&-9\\
    0&12&9
\end{bmatrix}\ ,\quad L=\begin{bmatrix}
    1&0&0\\
    2&1&0\\
    -2&\_&1
\end{bmatrix} 
$$
The last row operation is $r_3+2r_2$. Since $A$ in now upper triangular, we write $U$
$$
U=\begin{bmatrix}
    1&5&3\\
    0&-6&-9\\
    0&0&-9
\end{bmatrix}\ ,\quad L=\begin{bmatrix}
    1&0&0\\
    2&1&0\\
    -2&-2&1
\end{bmatrix}
$$
\bigskip
Firs use forward substitution to solve $Ly=b$:
$$
\begin{bmatrix}
    1&0&0\\
    2&1&0\\
    -2&-2&1
\end{bmatrix}
\begin{bmatrix} y_1\\y_2\\y_3 \end{bmatrix}
=\begin{bmatrix} 1\\-1\\-5 \end{bmatrix}
$$
\begin{align*}
y_1 &= 1\\
2y_1 + y_2 = -1 \implies y_2 &= -1-2(1) = -3\\
-2y_1 - 2y_2 + y_3 = -5 \implies y_3 &= -5 + 2(1) + 2(-3) = -9
\end{align*}
So $\mathbf{y}=\begin{bmatrix} 1\\-3\\-9 \end{bmatrix}$.

\bigskip
Then use backward substitution to solve $Ux=y$:
$$
\begin{bmatrix}
    1&5&3\\
    0&-6&-9\\
    0&0&-9
\end{bmatrix}
\begin{bmatrix} x_1\\x_2\\x_3 \end{bmatrix}
=\begin{bmatrix} 1\\-3\\-9 \end{bmatrix}
$$
\begin{align*}
-9x_3 = -9 \implies x_3 &= 1\\
-6x_2 - 9x_3 = -3 \implies x_2 &= \frac{-3+9(1)}{-6} = -1\\
x_1 + 5x_2 + 3x_3 = 1 \implies x_1 &= 1 - 5(-1) - 3(1) = 3
\end{align*}
$$
\mathbf{x}=\begin{bmatrix} 3\\-1\\1 \end{bmatrix}
$$
This matches the solution from question 1.

\bigskip

%%%%%%%%%%%%%%%%%%%%%%%%%%%%%%%%%%%%%%%%%%%%%%%%%%%%%%%%%%%%%%%%%%%%%%%%
\newpage
\section*{Question 3: QR decomposition (20 points)}

Perform QR decomposition for the matrix corresponding to the equations in \textbf{Question 1} using the Gram-Schmidt algorithm. Using the decomposition, solve for $\mathbf{x}$. Match your solution with that of \textbf{Question 1} and \textbf{Question 2}.

\textbf{Solution:}
\bigskip

%%%%%%%%%%%%%%%%%%%%%%%%%%%%%%%%%%%%%%%%%%%%%%%%%%%%%%%%%%%%%%%%%%%%%%%%
\newpage
\section*{Question 4: Gram-Schmidt process (20 points)}

Show that the residual vector $\mathbf{a}_i^\perp$ is orthogonal to $\mathbf{q}_1, \mathbf{q}_2, \dots, \mathbf{q}_{i-1}$ in the Gram-Schmidt process.

\textbf{Solution:}
\bigskip

%%%%%%%%%%%%%%%%%%%%%%%%%%%%%%%%%%%%%%%%%%%%%%%%%%%%%%%%%%%%%%%%%%%%%%%%
\newpage
\section*{Question 5: Rank deficient matrices (20 points)}

Compute the QR decomposition of the following two matrices using the Gram-Schmidt process and properties of the matrix $\mathbf{Q}$.
\begin{enumerate}
    \item[(i)] $\mathbf{A} = \begin{bmatrix}
        1 & 1 & -1\\
        1 & -1 & 1\\
        -1 & -1 & 1
    \end{bmatrix}$
    \item[(ii)] $\mathbf{B} = \begin{bmatrix}
        1 & -1 & 1\\
        -1 & 1 & -1\\
        -1 & 1 & -1
    \end{bmatrix}$
\end{enumerate}
Note: the matrices $\mathbf{Q}$ and $\mathbf{R}$ must be square matrices.

\textbf{Solution:}
\bigskip

\textbf{(i)}

\bigskip

\textbf{(ii)}

\end{document}
